% Copyright (c) 2026 Geovana Neves. Licensed under CC BY 4.0 (https://creativecommons.org/licenses/by/4.0/). See LICENSE-AND-AUTHORSHIP.md.
%
\begin{frame}{The physics behind the file, and where to read it}
\begin{itemize}\small
  \item \textbf{Coefficients} divide a force by the dynamic pressure and the
        reference area, and a moment by the area and a reference length; the
        moment is taken about the moment point [4].
  \item \textbf{The body axes} of the moment point and the rotor frames
        follow the right-hand rule of flight mechanics [4].
  \item \textbf{The advance ratio} \(J = V/(nD)\) and the propeller
        coefficients are the classical airscrew quantities [5, 6].
  \item \textbf{The actuator disc} is the momentum theory of a propeller: a
        pressure jump and a swirl across a disc of the rotor's radius [5].
  \item \textbf{The tip and helical Mach numbers} measure how far the blade
        tip is from the sonic limit of a linearised panel method [3, 7].
\end{itemize}
\end{frame}

\begin{frame}{References}
\begin{reflist}
\bibitem{g03r01} \docref{workspace-and-workflows}{The workspace and the workflow}
\bibitem{g03r02} \docref{getting-started}{Getting started}
\bibitem{g03r03} \docref{post-processing-definitions}{Post-processing definitions}
\bibitem{g03r04} B. Etkin, L. D. Reid, \emph{Dynamics of Flight: Stability and
      Control}, 3rd ed., Wiley, New York, 1996.
\bibitem{g03r05} H. Glauert, \emph{The Elements of Aerofoil and Airscrew Theory}, 2nd
      ed., Cambridge University Press, Cambridge, 1947.
\bibitem{g03r06} B. W. McCormick, \emph{Aerodynamics, Aeronautics, and Flight
      Mechanics}, 2nd ed., Wiley, New York, 1995.
\bibitem{g03r07} J. G. Leishman, \emph{Principles of Helicopter Aerodynamics}, 2nd
      ed., Cambridge University Press, Cambridge, 2006.
\end{reflist}
\end{frame}
